% 印刷用全展开封面：封底 + 书脊 + 正封，四周各 3 mm 出血
% 书脊宽度 \spinew（mm）由 排版/封面/做封面.sh 按页数算好后写入
\documentclass[tikz,border=0pt]{standalone}
\input{排版/封面/公共.tex}
\input{排版/build/封面/书脊宽度.tex}
\begin{document}
\begin{tikzpicture}[x=1mm,y=1mm]
  \pgfmathsetmacro{\S}{\spinew}
  \useasboundingbox (-3,-3) rectangle (343+\S,243);
  % ---------- 封底（x: 0–170） ----------
  \fill[cbg] (-3,-3) rectangle (170,243);
  \begin{scope}
    \clip (-3,-3) rectangle (170,243);
    \draw[cline!35, line width=0.15pt] (-3,-3) grid[step=10] (170,243);
  \end{scope}
  \node[anchor=north west, font=\heiti\bfseries\fontsize{17}{22}\selectfont, text=cink] at (18,222) {会翻跟头的机器人，是怎么造出来的？};
  \draw[cacc, line width=1.8pt] (18,211) -- (44,211);
  \node[anchor=north west, text width=134mm, font=\songti\fontsize{10.5}{19}\selectfont, text=cink, align=justify] at (18,205)
    {这本书以宇树的机器狗和人形机器人为主线，依据它的招股说明书、官方参数和开源代码，做了一次“纸上拆机”：从外壳、关节电机、减速器、电池、传感器，一路拆到平衡控制、强化学习和具身大模型；再用一手财务数据，讲清这门生意为什么能越卖越便宜、毛利率反而越来越高。最后两章写给想进入这一行的人：一张岗位地图，和一份不买真机也能完成的 90 天项目计划。};
  \node[anchor=north west, font=\heiti\bfseries\fontsize{12}{14}\selectfont, text=cink] at (18,150) {你会读到};
  \foreach \t [count=\i] in {十分钟读懂一张机器人参数表,机器人的后空翻和秧歌是怎么“练”出来的,用招股书算出来的单台成本和毛利,具身大模型今天走到了哪一步、卡在哪里,从自动驾驶、电机电控、互联网转行的接口在哪} {
    \fill[cacc] (19,142-\i*9+1.6) rectangle (21.6,142-\i*9+4.2);
    \node[anchor=west, font=\songti\fontsize{10.5}{13}\selectfont, text=cink] at (24,142-\i*9+2.9) {\t};
  }
  \node[anchor=north west, font=\heiti\bfseries\fontsize{12}{14}\selectfont, text=cink] at (18,82) {适合谁读};
  \node[anchor=north west, text width=134mm, font=\songti\fontsize{10.5}{17}\selectfont, text=cink] at (18,75)
    {看过机器人表演、想知道门道的普通读者；想转行具身智能的工程师和学生；关心机器人产业的从业者和投资者。};
  % 条码区
  \fill[white] (112,10) rectangle (160,40);
  \draw[cline, line width=0.4pt] (112,10) rectangle (160,40);
  \node[font=\heiti\fontsize{8}{10}\selectfont, text=cline, align=center] at (136,25) {ISBN 条码区\\（出版时放置）};
  \node[anchor=south west, font=\heiti\fontsize{8.5}{10}\selectfont, text=cink!80, align=left] at (18,30) {上架建议：科普读物 / 机器人 / 人工智能};
  \node[anchor=south west, font=\heiti\fontsize{8.5}{10}\selectfont, text=cink!80] at (18,23) {定价：（待定）};
  \node[anchor=south west, text width=88mm, font=\heiti\fontsize{7}{9}\selectfont, text=cink!65, align=left] at (18,10) {本书为独立写作，与宇树科技股份有限公司无隶属、授权或赞助关系。书中商标归各自所有者所有。};
  % ---------- 正封（x: 170+S – 340+S） ----------
  \begin{scope}[shift={(170+\S,0)}]
    \input{排版/封面/正面.tex}
  \end{scope}
  % ---------- 书脊（x: 170–170+S），最后画，盖在正封出血上面 ----------
  \fill[cink] (170,-3) rectangle (170+\S,243);
  \fill[cacc] (170,228) rectangle (170+\S,234);
  \node[anchor=north, text width=\S mm, align=center, font=\heitiblack\fontsize{16}{19}\selectfont, text=white] at (170+\S/2,220) {拆\\解\\宇\\树};
  \node[anchor=north, text width=\S mm, align=center, font=\heiti\fontsize{7}{8.4}\selectfont, text=white!85] at (170+\S/2,140) {从\\一\\颗\\电\\机\\到\\会\\翻\\跟\\头\\的\\机\\器\\人};
  \node[anchor=south, text width=\S mm, align=center, font=\heiti\bfseries\fontsize{10}{12}\selectfont, text=white] at (170+\S/2,14) {甜\\菜\\著};
\end{tikzpicture}
\end{document}
